\documentclass[11pt]{article}

\usepackage[utf8]{inputenc}
\usepackage[T1]{fontenc}
\usepackage{lmodern}
\usepackage{geometry}
\usepackage{amsmath,amssymb,amsfonts,amsthm,mathtools}
\usepackage{enumitem}
\usepackage{microtype}
\usepackage{hyperref}

\geometry{margin=1in}
\hypersetup{colorlinks=true, linkcolor=black, urlcolor=blue, citecolor=black}
\setlist[enumerate]{topsep=0.35em,itemsep=0.3em}
\setlist[itemize]{topsep=0.3em,itemsep=0.25em}

\newcommand{\Aether}{\AE{}ther}
\newcommand{\Aflow}{\AE{}ther-flow}
\newcommand{\TheoryProgram}{\Aether{} / \Aflow{} framework}
\newcommand{\Lang}{\mathcal{L}_{0}}
\newcommand{\LangPlus}{\mathcal{L}_{0}^{+H}}
\newcommand{\PiH}{\Pi_H}
\newcommand{\AdH}{\mathsf{Ad}_{H}}
\newcommand{\BlindH}{\mathsf{Blind}_{H}}
\newcommand{\GenH}{\mathsf{Gen}_{H}}
\newcommand{\EqSrc}{\mathsf{Eq}_{\mathrm{src}}}
\newcommand{\RetainH}{\mathsf{Retain}_{H}}
\newcommand{\NoTargetH}{\mathsf{NoTarget}_{H}}
\newcommand{\NoTargetSyn}{\mathsf{NoTarget}_{H}^{\mathrm{syn}}}
\newcommand{\AuditH}{\mathsf{Audit}_{H}}
\newcommand{\Tok}{\mathsf{Tok}}
\newcommand{\SrcTok}{\mathsf{Tok}_{0}}
\newcommand{\ForbidT}{\mathsf{Forbid}_{T}}
\newcommand{\ProxyT}{\mathsf{Proxy}_{T}}
\newcommand{\TraceH}{\mathsf{Trace}_{H}}
\newcommand{\Annot}{\mathsf{Ann}}
\newcommand{\Influence}{\mathsf{Infl}}
\newcommand{\Pass}{\mathsf{Pass}}
\newcommand{\Fail}{\mathsf{Fail}}
\newcommand{\Block}{\mathsf{Block}}
\newcommand{\Unres}{\mathsf{Unresolved}}

\newtheorem{selection}{Repair-target selection}
\newtheorem{target}{Theorem target}
\newtheorem{condition}{Failure condition}
\newtheorem{constraint}{Control constraint}
\newtheorem{verdict}{Local verdict}
\newtheorem{nextstep}{Next step}

\title{\textbf{Localization Source-Basis Axiom Selector-Domain NoTarget Source-Syntax Repair-Target Packet}\\[0.35em]
\large Least-Broad Mechanization Target for Target-Exclusion Syntax}
\author{Ontology Formalizer AgentJob AJ-RT-20260613-017-001}
\date{June 13, 2026}

\begin{document}

\maketitle

\begin{abstract}
This draft/control artifact selects \(\NoTargetH\) as the least broad next
repair target after the conservativity and benchmark-independence refutation
of the noncanonical ledger-generation ontology-expansion proposal. The
packet does not prove \(\NoTargetH\), adopt \(\LangPlus\), edit canonical
ontology sources, authorize candidate reconstruction, promote benchmark
status, or request Gate Chair review. It specifies a source-syntax theorem
target, \(\NoTargetSyn\), with explicit pass and failure conditions. The
result is a controlled-pause repair target: target exclusion must become a
mechanical syntax-and-provenance test before \(\EqSrc\), \(\RetainH\), or
\(\GenH\) can be evaluated without hidden target-proxy dependence.
\end{abstract}

\section{Scope}

This artifact is a draft/control output for task \texttt{RT-20260613-017}.
It responds to the prior finding that admissible variation of provenance
packets, admissibility rules, blinding records, source equivalence,
retention schemas, and target-exclusion grammars can change generated
families or hide adverse alternatives. The present packet performs no
candidate reconstruction and makes no canonical ontology change.

The repair question is:
\[
    \text{Which one of } \NoTargetH,\ \EqSrc,\ \RetainH,\ \GenH
    \text{ should be mechanized first?}
\]

\begin{selection}[Least-broad repair target]
\(\NoTargetH\) is the least broad first repair target because it is a
guard on the admissible syntax, annotations, relevance rules, and influence
relations used by the other three burdens. If target-proxy terms can enter
the generation trace, then a later equivalence relation, retention schema,
or generator theorem can inherit target selection while appearing
source-side.
\end{selection}

\begin{constraint}[Nonpromotion]
Selecting \(\NoTargetH\) as a repair target does not prove that
\(\NoTargetH\) holds. It only specifies the next theorem target and the
conditions under which the branch remains controlled-pause evidence.
\end{constraint}

\section{Source-Syntax Object}

Let \(\TraceH\) denote the proposed ledger-generation trace object for a
candidate packet. A trace includes:
\[
    \TraceH =
    (E,\ R,\ Q,\ A,\ B,\ S),
\]
where \(E\) is the finite expression set used in declarations and rules,
\(R\) is the relevance or admissibility-rule set, \(Q\) is the quotient or
equivalence grammar, \(A\) is the annotation and provenance record set,
\(B\) is the blinding record set, and \(S\) is the status and retention
schema.

The repair target requires the following source-side inventories:
\[
    \SrcTok,\quad \ForbidT,\quad \ProxyT,\quad \Annot,\quad \Influence.
\]
Here \(\SrcTok\) is the registered source-token inventory available to the
packet, \(\ForbidT\) is the declared target-token exclusion inventory,
\(\ProxyT\) is a declared proxy-risk inventory, \(\Annot\) is the
provenance annotation grammar, and \(\Influence\) is a syntactic influence
relation that records which tokens, annotations, labels, and relevance
choices can affect admission, equivalence, retention, ranking, or status.

\section{Theorem Target}

\begin{target}[NoTarget source-syntax theorem target]
The repair target is to prove or fail the following theorem in a future
bounded packet:

For every trace \(\TraceH=(E,R,Q,A,B,S)\) admitted by the noncanonical
proposal, \(\NoTargetSyn(\TraceH)=\Pass\) if and only if all of the
following are source-mechanical:
\begin{enumerate}
    \item every token occurring in \(E,R,Q,A,B,S\) is either in
    \(\SrcTok\) or is introduced by an explicit source-side declaration;
    \item no token in \(\ForbidT\) occurs in \(E,R,Q,A,B,S\);
    \item no element of \(\ProxyT\) has an \(\Influence\)-edge into
    admission, equivalence, retention, ranking, status, or family selection;
    \item every relevance or exclusion rule has a provenance annotation in
    \(\Annot\) that can be checked without benchmark outcome inspection;
    \item unclassified tokens, unannotated relevance rules, and proxy-risk
    influence edges return \(\Block\), not \(\Pass\); and
    \item semantic target-equivalence judgment is not used to classify any
    source expression as allowed.
\end{enumerate}
\end{target}

The theorem target is intentionally syntactic. It is not a claim that the
semantic concept of target independence has been solved. It is a proposed
mechanical guard that forces undecidable or semantically target-dependent
cases into \(\Block\) rather than silently passing.

\section{Failure Conditions}

\begin{condition}[Incomplete source-token inventory]
If \(\SrcTok\) is not complete for the packet under review, then
\(\NoTargetSyn\) cannot return \(\Pass\). The result is \(\Block\) until the
token inventory is registered or the packet is narrowed.
\end{condition}

\begin{condition}[Forbidden target-token occurrence]
If a token in \(\ForbidT\) occurs in \(E,R,Q,A,B,S\), the trace fails:
\[
    \NoTargetSyn(\TraceH)=\Fail.
\]
This covers direct target metric, target proper-time, target locality,
target observer protocol, target benchmark label, target-favorable rank, and
benchmark-success terms when they are declared in \(\ForbidT\).
\end{condition}

\begin{condition}[Proxy influence]
If a proxy-risk label or construct in \(\ProxyT\) influences admission,
equivalence, retention, ranking, status, or family selection, then the trace
does not pass. If the influence relation is explicit and source-mechanical,
the result is \(\Fail\). If the influence relation is underdetermined, the
result is \(\Block\).
\end{condition}

\begin{condition}[Semantic target-equivalence dependence]
If deciding whether a token, label, rule, or relevance choice is target-like
requires semantic comparison to target metric behavior or benchmark success,
then the proposed syntactic repair cannot discharge the guard. The branch
must remain controlled-pause evidence.
\end{condition}

\begin{condition}[Unannotated relevance or exclusion rules]
If relevance rules, exclusion rules, quotient choices, retention choices, or
status-row choices lack source-side provenance annotations, then
\(\NoTargetSyn\) returns \(\Block\). A hidden relevance choice is not allowed
to pass as source syntax.
\end{condition}

\section{Why EqSrc RetainH and GenH Remain Deferred}

\(\EqSrc\), \(\RetainH\), and \(\GenH\) are broader repair targets because
each can depend on the admissible trace language. Without \(\NoTargetSyn\):

\begin{itemize}
    \item \(\EqSrc\) can collapse or distinguish ledgers using hidden
    target-proxy criteria.
    \item \(\RetainH\) can retain adverse and favorable ledgers at unequal
    evidential grain through target-proxy relevance choices.
    \item \(\GenH\) can generate a favorable family from a trace that was
    source-clean only by assertion.
\end{itemize}

Therefore, the packet does not attempt to prove these burdens. It constrains
their future inputs by making the target-exclusion guard explicit first.

\section{Repair Ledger}

\begin{center}
\begin{tabular}{p{0.21\linewidth}p{0.34\linewidth}p{0.31\linewidth}}
\textbf{Surface} & \textbf{Repair target} & \textbf{Controlled-pause trigger} \\
\hline
\(\SrcTok\) & registered source-token inventory & incomplete inventory or undeclared extension \\
\(\ForbidT\) & target-token exclusion list & forbidden token occurs in trace \\
\(\ProxyT\) & proxy-risk labels and constructs & proxy influences admission ranking status or retention \\
\(\Annot\) & provenance grammar & relevance or exclusion rule lacks source annotation \\
\(\Influence\) & syntactic influence relation & influence cannot be checked source-locally \\
\(\NoTargetSyn\) & mechanical pass fail block predicate & semantic target-equivalence judgment is required \\
\end{tabular}
\end{center}

\section{Local Verdict}

\begin{verdict}[Repair target specified but not discharged]
This packet specifies \(\NoTargetSyn\) as a narrow theorem target with
failure conditions. It does not prove \(\NoTargetH\), does not adopt
\(\LangPlus\), and does not establish conservativity or benchmark
independence. The branch remains controlled until a later packet proves the
syntactic guard or records a failure.
\end{verdict}

\begin{verdict}[Candidate-support use remains blocked]
Candidate reconstruction, canonical ontology edit, benchmark promotion, Gate
Chair review, and completed-derivation language remain blocked. The result
is useful because it converts one broad vulnerability into a testable
source-syntax target.
\end{verdict}

\section{Next Role}

\begin{nextstep}[Smuggling audit before theorem attempt]
The logical next role is Smuggling Auditor. The auditor should test this
repair target for hidden target import, proxy-selection leakage, annotation
laundering, influence-relation laundering, and authority laundering before
any theorem attempt treats \(\NoTargetSyn\) as usable source-side guard
machinery.
\end{nextstep}

\section*{References}

Ricciardi, A. S. (2026, June 13). \emph{Localization source-basis axiom
selector-domain ledger-generation ontology-expansion proposal
conservativity and benchmark-independence refutation: Variation test for
provenance, admissibility, blinding, source equivalence, retention, and
target exclusion} [Unpublished manuscript]. The Aether-Flow Interpretation
of Relativity Research Project,
\texttt{research\_control/tasks/RT-20260613-016/artifacts/32\_LOCALIZATION\_SOURCE\_BASIS\_AXIOM\_SELECTOR\_DOMAIN\_LEDGER\_GENERATION\_ONTOLOGY\_EXPANSION\_PROPOSAL\_CONSERVATIVITY\_BENCHMARK\_INDEPENDENCE\_REFUTATION.tex}.

Ricciardi, A. S. (2026, June 13). \emph{Localization source-basis axiom
selector-domain ledger-generation ontology-expansion proposal: Noncanonical
source-language authority candidate for family-first ledger provenance}
[Unpublished manuscript]. The Aether-Flow Interpretation of Relativity
Research Project,
\texttt{research\_control/tasks/RT-20260613-010/artifacts/30\_LOCALIZATION\_SOURCE\_BASIS\_AXIOM\_SELECTOR\_DOMAIN\_LEDGER\_GENERATION\_ONTOLOGY\_EXPANSION\_PROPOSAL.tex}.

\end{document}
